\section{Discussion}
Hybrid-Evidence MIL distinguished MALT lymphoma from SS in the evaluated salivary gland material using slide-label training and fixed UNI2-h features. The internally selected checkpoint achieved internal validation AUROC 0.9599 and external AUROC 0.8721. Combining patch-level summaries provided a practical way to evaluate this diagnostic distinction without exhaustive regional annotations. These results motivate further assessment of the classification signal in cohorts with verified patient linkage and specimen metadata.

Hybrid had the highest internal AUROC among the evaluated individual models, while score averaging produced the highest ensemble AUROC. This descriptive ranking does not establish superiority over ABMIL: the paired AUROC-difference intervals included zero in both cohorts. Mean-feature logistic regression also retained external discrimination, with lower sensitivity at threshold 0.5. Representation quality, optimization and aggregation therefore need to be considered together, as emphasized in recent pathology benchmarks and nnMIL \citep{ref25,ref26}. The complete implementation comparisons support the reported performance profile, while the supplementary configuration analyses do not establish an isolated benefit of the learned gate.

External composition materially affected the clinical interpretation. Only 16 SS slides informed specificity, and MALT represented 90.48\% of the evaluated set. The reported accuracy of 91.07\% was therefore close to a constant-positive classifier. Excluding IHC/special-stain archive tags left AUROC similar but reduced sensitivity from 93.42\% to 75.61\%. This does not prove that staining caused the difference: source, diagnosis, stain and specimen characteristics were not separated by the available design. It does show why a single mixed-material sensitivity cannot stand for verified H\&E performance. The low external negative predictive value also prevents a rule-out interpretation. Class-weighted training and the absence of calibration fitting further limit interpretation of the sigmoid score as an individual probability.

The principal design limitation is the absence of an untouched internal test set. The same 49 slides selected epochs and supported model, ensemble and threshold comparisons. Fixed-score bootstrap intervals quantify uncertainty conditional on those predictions, not uncertainty from model selection. The split was performed by slide, and hospital-held patient linkage is needed to verify that patients did not span partitions or contribute correlated observations. Later analyses also reused an external collection whose historical outcomes were already known. Locking candidate selection before new external inference constrains that development exercise, but does not turn the collection into a fresh confirmatory cohort.

The region displays make selected image content traceable, while leaving biological interpretation open. Prototype margins are internal similarity differences, and coordinate-group pooling does not encode explicit interregional spatial relationships. The largest gate coefficient need not identify the feature with the largest causal contribution. Regional pathological annotations would be required to connect these outputs with lymphoepithelial lesions, glandular architecture or artifacts. Work such as SMMILe evaluates regional predictions against spatial evidence \citep{ref27}; a heatmap alone supplies a different level of evidence. Similarly, reader aggregates cannot establish an assisted-review benefit, which requires a defined reading protocol and individual paired outcomes \citep{ref21,ref22}.

The next clinical evaluation should first reconcile patient identifiers, diagnoses, source/stain metadata and eligibility decisions, then use patient-disjoint development with a frozen model and operating point assessed in an untouched, appropriately specified cohort. More SS controls are needed for precise specificity estimates. A protocolized reader study and annotation of selected regions would address clinical contribution and morphological relevance. These steps would test the application suggested by the present discrimination results without assuming that model development has already established diagnostic utility.
